%! AP Statistics — Unit 4, Lesson 4.12 — Warm-Up KEY
\documentclass[10pt]{article}
\usepackage{apstats-article}
\usepackage{apstats-key}

\begin{document}

\pageheader{Unit 4, Lesson 4.12}{Warm-Up --- Answer Key}
\begin{tabularx}{\linewidth}{X X X}
Name:\;\ans{Key} & Date:\; & Period:\;
\end{tabularx}

\vspace{0.3in}

{\large\bfseries\color{navy} Part 1 \quad Spiral Review}

\vspace{0.12in}

\begin{enumerate}

    \item $X \sim B(8, 0.25)$.
    \begin{enumerate}[label=(\alph*), itemsep=8pt]
        \item \ans{$\mu_X = 8(0.25) = 2$; \quad $\sigma_X = \sqrt{8(0.25)(0.75)} = \sqrt{1.5} \approx 1.22$}
        \item \ans{$P(X=2) = \binom{8}{2}(0.25)^2(0.75)^6 = 28(0.0625)(0.1779) \approx 0.3115$}
        \item \ans{$P(X \ge 3) = 1 - \texttt{binomcdf(8, 0.25, 2)} \approx 1 - 0.6785 = 0.3215$}
    \end{enumerate}

    \vspace{0.18in}

    \item
    \begin{enumerate}[label=(\alph*), itemsep=8pt]
        \item \ans{Geometric} --- \ans{$n$ is not fixed; rolling continues until the first 6.}
        \item \ans{Binomial} --- \ans{$n = 10$ is fixed; $X$ counts the number of 6s.}
        \item \ans{Neither} --- \ans{Sampling without replacement from a small deck violates the independence condition; neither binomial nor geometric applies strictly.}
    \end{enumerate}

    \vspace{0.18in}

    \item Basketball free throws, $p = 0.70$.
    \begin{enumerate}[label=(\alph*), itemsep=6pt]
        \item \ans{The only outcome where the first success is on the 3rd shot is: Miss, Miss, Make (MMK).}
        \item \ans{$P(\text{MMK}) = (0.30)(0.30)(0.70) = 0.063$}
    \end{enumerate}

\end{enumerate}

\vspace{0.2in}

{\large\bfseries\color{navy} Part 2 \quad Looking Ahead}

\vspace{0.12in}

\noindent Problems 2 and 3 preview the key contrast for today: in a \textbf{geometric}
setting, trials continue until the first success --- $n$ is not fixed.

\vspace{0.14in}

\begin{tcolorbox}[colback=greenbg, colframe=greenacc, boxrule=0.8pt, arc=2mm,
    left=4mm, right=4mm, top=3mm, bottom=3mm]
    \small
    \begin{itemize}[leftmargin=1.3em, itemsep=8pt]
  \item Check BITS (not BINS) for a geometric setting: the T stands for
        ``trials until the first success,'' replacing N (fixed number).
  \item $P(X = k) = (1-p)^{k-1}p$ --- multiply $k-1$ failure probabilities,
        then one success.
  \item $\mu_X = 1/p$ --- larger $p$ means fewer expected trials.
    \end{itemize}
\end{tcolorbox}

\vspace{0.18in}

\noindent\textcolor{slate}{\small Student responses will vary.}

\end{document}
